\documentclass[11pt, a4paper, oneside]{article}
\usepackage[ngerman]{babel}
\usepackage{worksheet}

\begin{document}
	\author{L. Bung}
	\title{Projekt: Spieleentwicklung}
	\subject{Informatik}
	\class{FTE1}
	\maketitle
	
	Das Ziel dieses Softwareprojekts ist es, ein funktionsfähiges Computerspiel zu entwickeln.
	Es soll sich dabei um ein textbasiertes Spiel handeln, welches über geschriebene Nutzereingaben gesteuert wird.
	Zusätzlich sollen der Ablauf des Spiels in eine Logdatei geschrieben werden.
	
	
	\section{Anforderungen an das Projekt}
	
	Inhaltlich gibt es keine Anforderungen an das Projekt.
	Sie können selbst entscheiden, um was es in dem Spiel geht -- beispielsweise ein Tic-Tac-Toe-Spiel, Mensch-ärgere-dich-nicht oder ein textbasiertes Rollenspiel.
	
	Das Projekt soll folgende Kriterien erfüllen:
	
	\begin{itemize}
		\item Sie müssen jedes der folgenden Konzepte in Ihrem Programm verwenden:
		\begin{itemize}
			\item Verzweigung (mindestens ein Block mit \texttt{if} und \texttt{else})
			\item Schleife (mindestens eine \texttt{while}-Schleife und eine \texttt{for}-Schleife)
			\item Liste
		\end{itemize}
		\item Jede logische Einheit des Programms muss als eigenständige Funktion geschrieben werden.
		\item Das Programm ist in angemessene Module aufzuteilen, zum Beispiel \texttt{main.py}, \texttt{input.py} und \texttt{logic.py}.
		\item Der Spielablauf soll in eine Logdatei geschrieben werden.
	\end{itemize}
	
	\section{Bewertungskriterien}
	
	Die Bewertung des Projekts gliedert sich in mehrere Bestandteile:
	
	\begin{enumerate}
		\item Das \textbf{Python-Programm} an sich
		\item Eine schriftliche \textbf{Dokumentation} des Projekts
		\item \textbf{Sonstige Kriterien} wie z. B. Gesamteindruck, Kreativität usw.
	\end{enumerate}
	
	\subsection{Python-Programm}
	
	Neben den oben genannten Kriterien spielen weitere Kriterien eine Rolle, wie zum Beispiel die folgenden:
	
	\begin{itemize}
		\item \textbf{Kommentare}: Schreiben Sie ausführliche Kommentare.
		Immer dann, wenn nicht auf den ersten Blick erkennbar ist, was ein Codeabschnitt tut, sollten Sie dies erläutern.
		Funktionen sollten immer einen Kommentar haben, der beschreibt, was die Funktion tut.
		\item \textbf{Lesbarkeit}: Achten Sie darauf, den Code so lesbar wie möglich zu schreiben.
		Dazu gehört z. B. eine sinnvolle Benennung der Variablen usw.
		\item \textbf{Struktur}: Strukturieren Sie Ihren Code logisch.
		Dazu gehört die Aufteilung in Funktionen und Module, aber auch innerhalb der einzelnen Abschnitte ist eine sinnvolle Strukturierung wichtig.
	\end{itemize}
	
	\subsection{Dokumentation}
	
	Die Dokumentation dient einerseits dazu, um einen Überblick über Ihr Projekt zu liefern, ohne in den Code zu schauen.
	Andererseits hilft das Schreiben der Dokumentation Ihnen auch, damit Sie im Nachhinein den Entwicklungsprozess und Ihre Gedanken nachvollziehen zu können.
	
	In der Dokumentation sollten folgende Dinge vorhanden sein:
	
	\begin{itemize}
		\item Sie sollten den Inhalt und die Idee Ihres Projekts beschreiben.
		\item Der zeitliche Ablauf Ihrer Entwicklung sollte erkennbar werden.
		\item Überlegungen zu Struktur und Aufbau Ihres Programms sollten erklärt werden.
		\item Die Dokumentation sollte einen Abschnitt beinhalten, in dem erklärt wird, wie das Programm auszuführen ist.
	\end{itemize}
	
	Bitte geben Sie Ihre Dokumentation im PDF-Format ab, nicht als \texttt{.docx}-Datei o. Ä.
	
	\section{Abgabe}
	
	\begin{itemize}
		\item Die Abgabe erfolgt als Upload via Moodle.
		\item Bitte beachten Sie die dort angegebene Deadline!
		\item Eine spätere Abgabe ist nur in Sonderfällen nach Rücksprache und vorheriger Genehmigung möglich.
		\item Sie können mehrere Dateien mit einer Gesamtgröße von maximal 20 MB abgeben.
		\item Wenn Sie mehrere Dateien auf einmal abgeben, komprimieren Sie diese bitte als ZIP-Archiv.
		\item Sollten Sie in Ihrer Dokumentation bereits den Python-Code anhängen, bitte ich dennoch um zusätzliche Abgabe der Codedateien.
		\item Das Programm muss im Schulnetz unter der Linuxinstallation ausführbar sein.
	\end{itemize}
	
	\section{Tipps}
	
	\begin{itemize}
		\item Planen Sie Ihr Projekt nicht zu groß.
		In der Regel braucht man in der Softwareentwicklung etwa 1,5- bis 2-fach so viel Zeit wie zuvor absehbar.
		\item Arbeiten Sie kontinuierlich und versuchen Sie nicht, in einer Nachtschicht das komplette Programm zu schreiben.
		\item Kümmern Sie sich zuerst um die Logik Ihres Spiels, bevor Sie sich mit Optimierungen wie z. B. einer schön formatierten Nutzereingabe beschäftigen:
		
		``Premature optimization is the root of all evil'' -- Donald Knuth\footnote{\url{https://en.wikipedia.org/wiki/Program_optimization\#When_to_optimize}}
		\item Wenn eine Python-Datei so lang wird, dass unübersichtlich zu lesen wird, ist das ein relativ sicheres Zeichen dafür, dass Sie mehrere Funktionen und Module dafür verwenden sollten.
		\item Die Dokumentation ist das ``Aushängeschild'' Ihres Programms.
		Eine gut strukturierte, formulierte und vollständige Dokumentation zahlt sich aus.
		\item Schreiben Sie die Dokumentation während des Entwicklungsprozesses und nicht erst ganz am Ende.
		Zum Schluss können Sie diese nochmal überarbeiten.
	\end{itemize}
\end{document}